\begin{mistake}{2026-07-16}{03}

\begin{problemcontent}
设 \(a_n = \int_0^{\frac{\pi}{4}} \tan^n x \mathrm{d}x\)，证明：\(\frac{1}{2(n+1)} < a_n < \frac{1}{2(n-1)} \ (n \ge 2)\)。
\end{problemcontent}

\begin{solutioncontent}
【法一：递推与放缩】
计算 \(a_n + a_{n+2}\)：
\[ a_n + a_{n+2} = \int_0^{\frac{\pi}{4}} \tan^n x (1 + \tan^2 x) \mathrm{d}x = \int_0^{\frac{\pi}{4}} \tan^n x \mathrm{d}(\tan x) = \frac{1}{n+1} \]
同理，\(a_{n-2} + a_n = \frac{1}{n-1}\)。
在 \((0, \frac{\pi}{4})\) 上，\(0 < \tan x < 1\)，故指数越大值越小，即 \(\tan^{n+2} x < \tan^n x < \tan^{n-2} x\)。
两边积分得：\(a_{n+2} < a_n < a_{n-2}\)。
利用单调性放缩递推式：
\[ \frac{1}{n+1} = a_n + a_{n+2} < a_n + a_n = 2a_n \implies a_n > \frac{1}{2(n+1)} \]
\[ \frac{1}{n-1} = a_{n-2} + a_n > a_n + a_n = 2a_n \implies a_n < \frac{1}{2(n-1)} \]

【法二：换元与代数放缩】
令 \(t = \tan x\)，由于 \(x \in (0, \frac{\pi}{4})\)，则 \(t \in (0, 1)\)，\(\mathrm{d}x = \frac{1}{1+t^2} \mathrm{d}t\)。
代入原积分：
\[ a_n = \int_0^1 \frac{t^n}{1+t^2} \mathrm{d}t \]
在 \(t \in (0, 1)\) 内，有 \(t^2 < 1\)，因此 \(1+t^2 < 2\)，且 \(2t^2 < 1+t^2\)。
于是被积函数满足：
\[ \frac{1}{2} t^n < \frac{t^n}{1+t^2} < \frac{t^n}{2t^2} = \frac{1}{2} t^{n-2} \]
在 \([0, 1]\) 上对不等式两边积分，即得：
\[ \frac{1}{2(n+1)} < a_n < \frac{1}{2(n-1)} \]
\end{solutioncontent}

\mistakeknowledge{
\begin{itemize}
  \item \textbf{核心公式/定理}：\(\int \tan^n x \sec^2 x \mathrm{d}x = \frac{1}{n+1}\tan^{n+1}x\)；换元法 \(\mathrm{d}(\arctan t) = \frac{1}{1+t^2}\mathrm{d}t\)。
  \item \textbf{易错点}：遇到 \(\tan^n x\) 定积分，首选“加两阶”凑微分；或作 \(t=\tan x\) 转化为有理函数积分，避免陷入分部积分。
  \item \textbf{关联知识}：定积分保号性；有理分式在 \(0 < t < 1\) 区间内的不等式放缩。
\end{itemize}}

\end{mistake}
